\chapter{Guidance design}
\label{ch:guide}

The previous chapter specified the scenario to be generated through a mask
$\pi$ and a set of member-specific target trajectories. We now address the
second component of the framework: how the generative forecast should be
steered toward these prescribed targets.

Because guidance is applied independently to each ensemble member, we suppress
the ensemble-member index throughout this chapter. Accordingly, $y_n^\star$
denotes the target associated with the member currently being generated.

(TODO: this can probably be specified better connecting to something previous instead of restating again)
We formulate the problem as \emph{target-driven guidance}. Rather than choosing
a guidance strength beforehand and observing the resulting change in the
forecast, the desired outcome is specified first. Guidance strength then
becomes a means of realizing this target. At each weather step, we use the
target discrepancy evaluated on the intermediate clean-state estimates to
construct a local guidance direction, while the magnitude of the intervention
is adapted separately to close the remaining target gap at the end of the
generative flow.

Several strategies can be considered for this target-realization problem. One
possibility is to adapt the intervention locally such that the target gap
follows a prescribed evolution during the flow. A more general alternative is
to optimize the guidance strengths at all flow steps jointly. In this work, we
instead prescribe the relative distribution of guidance over flow time and
adapt only a single scalar controlling its overall magnitude. Preliminary
investigations of the former approaches are discussed in
\Cref{app:target-realization-development}; the scalar-gain formulation retained
here provides the basis for the experiments in this thesis.

TODO: here we can state that it becomes clear from the experiments why this formulation doesn't demand for any optimization of the schedule under the DPS formulation.


\section{Target-driven guidance}
\label{sec:target-driven-guidance}

ArchesWeather generates a forecast trajectory autoregressively. Given the
atmospheric context available at weather step $n$, the model generates the next
state $x_{n+1}$. Guidance is applied within the residual flow used to generate
this state.

Although the guidance problem is solved separately at each weather step, its
effects propagate through weather time. Once a guided state
$x_{n+1}^{\mathrm{gui}}$ has been generated, it becomes part of the
conditioning context used to generate the following state. Consequently,
guidance applied at one weather step can influence all subsequent states of the
guided forecast trajectory.


\subsection{Guidance objective}
\label{sec:guidance-objective}

At weather step $n$, the specification procedure introduced in
\Cref{ch:specify} provides the target value $y_{n+1}^\star$ for the state being
generated. During the residual flow, ArchesWeather provides a clean-state
estimate $\hat{x}_{n+1,t}$ at each flow step $t$. We define its discrepancy
from the prescribed target as

\begin{equation}
    \xi_{n+1,t}
    :=
    \mathcal{M}_\pi
    \left(
        \hat{x}_{n+1,t}
    \right)
    -
    y_{n+1}^\star,
    \label{eq:guidance-gap}
\end{equation}

where $\mathcal{M}_\pi$ denotes the masked functional defined in
\Cref{ch:specify}. We refer to $\xi_{n+1,t}$ as the \emph{target gap} at flow
step $t$.

The corresponding guidance objective is the squared target gap,

\begin{equation}
    \mathcal{L}_{n+1,t}
    :=
    \xi_{n+1,t}^{\,2}
    =
    \left[
        \mathcal{M}_\pi
        \left(
            \hat{x}_{n+1,t}
        \right)
        -
        y_{n+1}^\star
    \right]^2.
    \label{eq:guidance-loss}
\end{equation}

Evaluating the objective on the clean-state estimate follows the local
clean-estimate guidance formulation introduced in
\Cref{sec:guidance-as-posterior-sampling}. It makes the conditioning objective
available throughout the generative flow and therefore provides a guidance
direction at every intermediate state.

An alternative would be to define the objective only on the final generated
state and differentiate it backwards through the complete flow. Such a
formulation is possible and corresponds to the control perspective discussed
in \Cref{sec:guidance-as-control}. The distinction is therefore not whether
the final state can be differentiated through, but where the intervention is
defined. Here, we retain the local posterior-style construction: the clean
estimate determines the guidance direction at each flow step, while the
resulting terminal target gap is used separately to determine how strongly
this direction must be applied.

For the experiments in this thesis, $\mathcal{M}_\pi$ is the area-weighted
regional mean surface temperature introduced in \Cref{ch:specify}. The
guidance construction itself is more general, however, and can be applied to
other differentiable functionals of the forecast state.


\subsection{Gradient-based flow guidance}
\label{sec:gradient-flow-guidance}

We index the discrete residual flow such that
$z_{n+1,0}$ denotes the initial noise and $t=1,\ldots,T$ indexes the successive
flow updates,

\begin{equation}
    z_{n+1,0}
    \xrightarrow{\;1\;}
    z_{n+1,1}
    \xrightarrow{\;2\;}
    \cdots
    \xrightarrow{\;T\;}
    z_{n+1,T}.
    \label{eq:guided-flow-indexing}
\end{equation}

At flow step $t$, the model evaluates the velocity field and clean-state
estimate from the current noisy state $z_{n+1,t-1}$. The local guidance
direction is obtained by differentiating the loss in
\Cref{eq:guidance-loss} with respect to this state,

\begin{equation}
    g_{n+1,t}
    :=
    \nabla_{z_{n+1,t-1}}
    \mathcal{L}_{n+1,t}.
    \label{eq:guidance-gradient}
\end{equation}

For the squared target-gap objective, this can be written as

\begin{equation}
    g_{n+1,t}
    =
    2\xi_{n+1,t}\,
    \nabla_{z_{n+1,t-1}}
    \mathcal{M}_\pi
    \left(
        \hat{x}_{n+1,t}
    \right).
    \label{eq:guidance-gradient-expanded}
\end{equation}

Although the loss is defined through the masked quantity
$\mathcal{M}_\pi(\hat{x}_{n+1,t})$, the gradient is propagated through the
mapping from the noisy residual state to the complete clean-state estimate.
The resulting direction therefore describes how the current noisy state can
be modified, through the learned generative model, to change the targeted
forecast quantity. It is not restricted to modifying only the entries selected
by the mask.

Following the loss-based guidance formulation introduced in the background
chapter, the velocity field is modified according to

\begin{equation}
    \tilde{u}_{n+1,t}
    =
    u_{n+1,t}
    -
    \lambda_{n,t} g_{n+1,t},
    \label{eq:target-driven-guided-velocity}
\end{equation}

where $\lambda_{n,t}\geq0$ controls the guidance magnitude at weather step $n$
and flow step $t$. The resulting Euler update is

\begin{equation}
    z_{n+1,t}
    =
    z_{n+1,t-1}
    +
    h_t
    \tilde{u}_{n+1,t},
    \qquad
    t=1,\ldots,T.
    \label{eq:guided-euler-update}
\end{equation}

The local guidance direction is therefore fully determined by the specified
target and the current generative state. What remains to be determined is the
sequence of guidance magnitudes $\lambda_{n,t}$ required to realize the
prescribed target. We address this design problem next.

\section{Guidance-strength design}
\label{sec:guidance-strength-design}

The background formulation in \Cref{sec:guidance-strength-schedules} decomposes
the guidance strength into an overall magnitude, a flow-time profile, and an
optional normalization factor. Applied independently at each weather step, this
gives

\begin{equation}
    \lambda_{n,t}
    =
    w_n \gamma_t \nu_{n,t},
    \label{eq:guidance-strength-decomposition}
\end{equation}

where $w_n$ controls the overall guidance magnitude used to generate
$x_{n+1}$, $\gamma_t$ determines its relative variation over flow time, and
$\nu_{n,t}$ represents an optional additional normalization.

In our method, we apply the raw guidance gradient without an additional
normalization and therefore set

\begin{equation}
    \nu_{n,t}=1.
\end{equation}

The guidance strength consequently reduces to

\begin{equation}
    \lambda_{n,t}
    =
    w_n\gamma_t.
    \label{eq:scalar-gain-guidance-strength}
\end{equation}

We prescribe $\gamma_t$ as a non-negative flow-time profile and normalize it
such that

\begin{equation}
    \max_{t=1,\ldots,T-1}\gamma_t = 1.
    \label{eq:guidance-profile-normalization}
\end{equation}

Only the relative variation of $\gamma_t$ over flow time is therefore encoded
by the profile, while its overall magnitude is carried entirely by $w_n$.
Conceptually, the two quantities address different aspects of the guidance
problem: $\gamma_t$ determines \emph{where} along the generative flow guidance
is applied, whereas $w_n$ determines \emph{how strongly} it is applied at the
current weather step. The profile families considered in this thesis are
introduced together with the experimental design.

Substituting \Cref{eq:scalar-gain-guidance-strength} into
\Cref{eq:target-driven-guided-velocity} gives the update used throughout the
main experiments,

\begin{equation}
    \tilde{u}_{n+1,t}
    =
    u_{n+1,t}
    -
    w_n\gamma_t g_{n+1,t}.
    \label{eq:scalar-gain-guided-velocity}
\end{equation}


\subsection{Alternative target-realization strategies}

The decomposition above leaves several possibilities for determining the
guidance strengths required to realize the prescribed target. During the
development of the method, we considered three increasingly constrained
formulations.

A first approach adapts the guidance magnitude locally at every flow step such
that the remaining target gap follows a prescribed trajectory toward zero.
This directly controls the evolution of the gap during generation, but requires
specifying how quickly the target should be approached within the internal
flow.

A second possibility treats the individual guidance strengths
$\lambda_{n,1},\ldots,\lambda_{n,T-1}$ as independent optimization variables
and determines them jointly from the final target discrepancy. This provides
substantially more flexibility, but replaces the scalar guidance-strength
choice by a higher-dimensional optimization problem.

The formulation retained in this thesis instead fixes their relative structure
through $\gamma_t$ and adapts only the scalar gain $w_n$. Preliminary
experiments showed that this restricted formulation was sufficient to realize
the prescribed targets while reducing the target-realization problem to a
single dimension. The additional flexibility of the former approaches was
therefore not required for the experiments considered here. Their formulations
and their role in the development of the final method are documented in
\Cref{app:target-realization-development}.


\subsection{Final flow step}

Guidance is applied during the first $T-1$ flow updates but omitted from the
final update. This convention follows directly from the clean-prediction
parameterization of ArchesWeather.

At the final flow step, the current residual state is
$z_{n+1,T-1}$ and the remaining noise level is $s_T$. From
\Cref{eq:arches-velocity}, the corresponding velocity is

\begin{equation}
    u_{n+1,T}
    =
    \frac{
        \hat{r}_{n+1,T}
        -
        z_{n+1,T-1}
    }{
        s_T
    }.
    \label{eq:final-flow-velocity}
\end{equation}

Because the final Euler step integrates the remaining noise level completely,
its step size is $h_T=s_T$. Without guidance,

\begin{align}
    z_{n+1,T}
    &=
    z_{n+1,T-1}
    +
    s_T u_{n+1,T}
    \\
    &=
    z_{n+1,T-1}
    +
    \hat{r}_{n+1,T}
    -
    z_{n+1,T-1}
    \\
    &=
    \hat{r}_{n+1,T}.
    \label{eq:final-flow-snap}
\end{align}

The final update therefore maps the noisy residual state directly onto the
current clean-residual prediction. A guidance correction applied at this step
would perturb the resulting residual without any subsequent model evaluation
in which the modified state could be processed. We consequently set

\begin{equation}
    \gamma_T
    =
    \lambda_{n,T}
    =
    0,
    \label{eq:no-final-guidance}
\end{equation}

such that $t=T-1$ is the final guided flow step. Omitting guidance at the final
noise level is also consistent with the precipitation-guidance procedure of
\textcite{ueyama2026guided}.


\section{Target realization by gain adaptation}
\label{sec:target-realization}

Once the flow-time profile $\gamma_t$ has been fixed, the only remaining
guidance parameter at weather step $n$ is the scalar gain $w_n$. Rather than
choosing this gain beforehand, we determine it from the target that was
specified for the state $x_{n+1}$.

For a fixed atmospheric context, flow-time profile, and initial noise
realization $z_{n+1,0}$, running the complete guided residual flow defines the
generated state as a function of the gain,

\begin{equation}
    x_{n+1}^{\mathrm{gui}}(w_n).
\end{equation}

We define the corresponding terminal target gap by

\begin{equation}
    \xi_{n+1}(w_n)
    :=
    \mathcal{M}_\pi
    \left(
        x_{n+1}^{\mathrm{gui}}(w_n)
    \right)
    -
    y_{n+1}^\star.
    \label{eq:terminal-target-gap}
\end{equation}

This quantity differs from the intermediate gap
$\xi_{n+1,t}$ in \Cref{eq:guidance-gap}. The latter is evaluated on a clean
estimate during the flow and determines the local guidance direction, whereas
$\xi_{n+1}(w_n)$ is evaluated on the final generated state and measures whether
the prescribed target was actually realized.

The target-realization problem is therefore the scalar root-finding problem

\begin{equation}
    \xi_{n+1}(w_n^\star)
    =
    0.
    \label{eq:target-realization-root}
\end{equation}

This formulation directly reflects the target-driven nature of the method:
the desired endpoint is fixed first, and the role of $w_n$ is solely to
determine the guidance magnitude required to reach it.


\subsection{Secant search}

Since evaluating $\xi_{n+1}(w_n)$ requires running the complete guided
residual flow, we seek its root without computing derivatives with respect to
the gain. We use the secant method, which constructs a local linear
approximation from two previously evaluated gains.

Given $w_n^{(j-1)}$ and $w_n^{(j)}$, the next candidate is

\begin{equation}
    w_n^{(j+1)}
    =
    w_n^{(j)}
    -
    \xi_{n+1}
    \left(
        w_n^{(j)}
    \right)
    \frac{
        w_n^{(j)}-w_n^{(j-1)}
    }{
        \xi_{n+1}
        \left(
            w_n^{(j)}
        \right)
        -
        \xi_{n+1}
        \left(
            w_n^{(j-1)}
        \right)
    }.
    \label{eq:secant-guidance-gain}
\end{equation}

The first evaluation uses $w_n=0$, corresponding to generation of the current
weather step without flow guidance, and a second evaluation uses a positive
initial gain. For later weather steps, the $w_n=0$ evaluation should not be
confused with the matched unguided trajectory used to define the target in
\Cref{ch:specify}. Previous guided states already form part of the
autoregressive conditioning context, so removing guidance from the current
residual flow does not in general recover the independently generated unguided
forecast.

All evaluations performed by the secant search use the same initial noise
realization $z_{n+1,0}$. Consequently, the mapping

\[
    w_n
    \longmapsto
    \xi_{n+1}(w_n)
\]

changes only through the guidance magnitude and not through stochastic
differences between repeated flow samples. Importantly, the local guidance
gradients $g_{n+1,t}$ are nevertheless recomputed during every evaluation,
since changing $w_n$ changes the sequence of noisy states and therefore the
clean-state estimates encountered along the flow.

The secant iteration is restricted to non-negative gains. If the target is
reached within the prescribed tolerance, the corresponding gain is retained.
If no root is obtained within the numerical safeguards or evaluation budget,
we instead select the evaluated value that leaves the smallest absolute
terminal gap,

\begin{equation}
    w_n^\star
    :=
    \arg\min_{w\in\mathcal{W}_n}
    \left|
        \xi_{n+1}(w)
    \right|,
    \label{eq:best-evaluated-gain}
\end{equation}

where $\mathcal{W}_n$ denotes the set of gains evaluated during the search.
The concrete initialization, convergence tolerance, step safeguards, and
maximum number of evaluations are specified together with the experimental
setup.

This procedure is repeated independently at every guided weather step. The
resulting sequence of gains

\[
    w_0^\star,\ldots,w_{H-1}^\star
\]

therefore adapts the overall guidance magnitude to the member-specific target
at each point along the prescribed trajectory. For weather states beyond the
guided horizon, $n+1>H$, no target is imposed and the autoregressive forecast
continues without guidance.